\documentclass[reqno]{amsart}
\usepackage{fullpage}
\usepackage{parskip}
\usepackage{graphicx}
\usepackage{float}
\usepackage{xcolor}
\newcommand{\1}{\mathbbm{1}}
% measures / sets of measures:
\newcommand{\mH}{\mathscr{H}}
\newcommand{\mL}{\mathscr{L}}
\newcommand{\mP}{\mathscr{P}}
\newcommand{\cH}{\mathcal{H}} % entropy
\newcommand{\cT}{\mathcal{T}} % topology
\newcommand{\cB}{\mathcal{B}} % borel
\newcommand{\W}{\mathcal{W}} % new metric W
\newcommand{\X}{\mathcal{X}} % space X
\newcommand{\Y}{\mathcal{Y}} % space Y
\newcommand{\expv}{\mathbb{E}} % exp value

\newcommand{\F}{\mathcal{F}} % space F
\newcommand{\M}{\mathcal{M}} % space of measures

\newcommand{\deriv}[2]{\frac{\de #1}{\de #2}}
\newcommand{\pderiv}[2]{\frac{\partial #1}{\partial #2}}
\newcommand{\sgn}[0]{\mathrm{sign}}

\newcommand{\nabladot}{\nabla \cdot} % discrete divergence operator
\newcommand{\A}{\mathcal{A}} % action
\newcommand*\de{\mathop{}\!\,\,\mathrm{d}}
\DeclareMathOperator*{\argmin}{arg\,min}
\DeclareMathOperator*{\argmax}{arg\,max}
\DeclareMathOperator*{\supp}{supp}
\DeclareMathOperator*{\id}{id}
\DeclareMathOperator{\imm}{Im}
\DeclareMathOperator{\lalpha}{\overline{\alpha}}
\DeclareMathOperator{\aconv}{\alpha \textit{x}+\overline{\alpha}\textit{y}}
\DeclareMathOperator{\bconv}{\beta \textit{x}+\overline{\beta}\textit{y}}

%%%%%%%%%%%%%%%%%%%% from the paper
\newcommand{\R}{\ensuremath{\mathbb{R}}}
\newcommand{\N}{\ensuremath{\mathbb{N}}}
\newcommand{\Z}{\ensuremath{\mathbb{Z}}}
\newcommand{\card}{\vert\mathcal{H}\vert} % |H|
\newcommand{\LL}{\mathcal L}
\newcommand{\Rad}{\widehat{\mathfrak R}_D}


\DeclareMathOperator{\Ric}{Ric}
\DeclareMathOperator{\vol}{vol}
\DeclareMathOperator{\Var}{Var}

\def\diam{\mathop{\mathrm{diam}}\nolimits}
\def\vol{\mathop{\mathrm{vol}}\nolimits}
\def\area{\mathop{\mathrm{area}}\nolimits}
\def\Hess{\mathop{\mathrm{Hess}}\nolimits}
\def\supp{\mathop{\mathrm{supp}}\nolimits}
\def\Ric{\mathop{\mathrm{Ric}}\nolimits}
\def\Id{\mathop{\mathrm{Id}}\nolimits}
\def\tr{\mathop{\mathrm{tr}}\nolimits}
\def\ac{\mathop{\mathrm{ac}}\nolimits}
\def\div{\mathop{\mathrm{div}}\nolimits}
\def\proj{\mathop{\mathrm{proj}}\nolimits}

%%%%%%%%%%%%%%%%%%%%%%%%%%%%%%%%%%%%

\theoremstyle{plain}
\newtheorem{definition}{Definition}[section]
\newtheorem{proposition}[definition]{Proposition}
\newtheorem{theorem}[definition]{Theorem}

\newtheorem{corollary}[definition]{Corollary}
\newtheorem{lemma}[definition]{Lemma}
\newtheorem{example}[definition]{Example}
\newtheorem{remark}[definition]{Remark}
\newtheorem{assumption}[definition]{Assumption}
\newtheorem{fact}[definition]{Fact}

\newenvironment{intuition}[1][Intuition]{\begin{proof}[#1]\renewcommand{\qed}{\relax}}{\topsep0pt\end{proof}}
\linespread{1.2}

\usepackage{hyperref}
\usepackage{enumitem}
\renewcommand{\labelenumi}{$(\roman{enumi})$}
\renewcommand{\labelitemi}{$\cdot$}
\usepackage{mathtools}
\usepackage{amssymb}
\usepackage{amsmath}

\newcommand{\loz}[1]{\textcolor{orange}{[LOZ:\ #1]}}
\newcommand{\mati}[1]{\textcolor{magenta}{#1}}
% \pie ? 

% add some color
\definecolor{benred}{HTML}{8B0000}
\definecolor{bluette}{RGB}{0,100,255} % bright bluish tone
\definecolor{darkorange}{RGB}{255, 140, 0}
\definecolor{pink}{RGB}{255, 130, 190}
\definecolor{acquagreen}{RGB}{0,230,255} % verde acqua

\title{Towards a structure-aware bound on the test error of decision trees}
\author{October 2026}
\begin{document}
\maketitle

\section*{Introduction}

{\bf Setting}

Given a sample space $\Omega$, let inputs and labels be generated by $X:\Omega \to \mathcal X = \{\bar x_1,\ldots,\bar x_N\}$ and $Y: \Omega \to \mathcal Y = \{-1, 1\}$. Let $p$ be a probability distribution on $\mathcal X\times\mathcal Y$, with
\[
 p_{j,y}=\Pr[X=\bar x_j,Y=y],\qquad
 p_{j,y}\ge0,\qquad
 \sum_{j=1}^N\sum_{y\in\mathcal Y}p_{j,y}=1.
\]
Let $D = {(X_i, Y_i)}_{i=1}^n$ be a dataset of $n$ indipendent copies of $(X,Y)$ with common distribution $p$. We indicate the realizations of $D$ with lower-case $x$ and $y$. 
%We want to predict $y$ from $x$ by learning a function $h:\mathcal X \to \mathcal Y$. 
Let $\mathcal H \subseteq \{h:\mathcal X \to \mathcal Y\}$ be a nonempty, finite set of classifiers. 
Given a loss function $\ell:\mathcal Y \times \mathcal Y \to \mathbb R^+$, define the population and empirical risks by
\[
    R(h) = \mathbb E_D(\ell(h(X), Y)) = \sum_{\bar x\in\mathcal X} \sum_{\bar y \in\{-1,1\}} \ell(h(\bar x), \bar y) p(\bar x, \bar y).
\]
and
\[
    \hat R_D (h) = \frac{1}{n}\sum_{i=1}^n \ell(h(X_i), Y_i).
\]
We select the classifier that minimizes the empirical risk (ERM), i.e.
\begin{equation*}\label{eq: ERM}
    \hat h_D \in \argmin_{h\in \mathcal H} \hat R_D (h)
\end{equation*}
which, importantly, is a random variable of the data $D$. Note that at least one such $\hat h_D$ exists, since $\cH$ is finite and nonempty.

Specifically, in this thesis 
\begin{itemize}
    \item $\ell$ is the counting loss $\ell(h(x), y) = \mathbf{1} (h(x)\neq y) = \frac{1 - h(x)y}{2}$;
    \item classifiers $h:\X\to\Y$ are  binary decision trees;
    \item given a depth $d$, $\mathcal H_d$ is the set of binary decision trees of depth at most $d$, and we consider a number of such depths. 
\end{itemize}
\bigskip 

{\bf Research question}

In machine learning, there exists a large number of algorithms, notably neural networks trained with gradient descent, which achieve great performances at prediction tasks on unseen data. A big question in learning theory is to understand the relationship between the behavior of such algorithms on the available data and their future behavior on unseen data (generalization). How can we predict how this future behavior will look like? 

A number of results have tried to answer this question. Recent works seem to suggest that the classical theory fails to explain generalization, in particular they fail to describe the behavior of the test error of learning algorithms as the comlpexity of the hypotheses set increases. This thesis seeks to understand whether this is true by testing theoretical results against empirical evidences. We study the setting of a finite hypotheses set of binary decision trees. We look for a function that reproduces qualitatively the test error as the maximum depth of the decision trees (i.e.\ the complexity of the hypotheses space) increases. The contribution of identifying and empirically validating such a function is two-fold. On the one hand, it explains the shape of the test error curve of decision trees; on the other hand, it may allow to implement the same framework and techniques to study the generalization performance of more complex algorithms, like neural networks trained with gradient descent. 

In particular, we study bounds on the expcted gap between the test and train error of the ERM, \[\expv_D\left[(R - \hat R_D)(\hat h_D)\right].\] 
Note that since $\hat h_D$ and $\hat R_D$ depend on the same dataset $D$, $\expv_D[\hat R_D(\hat h_D)]\neq R(\hat h_D)$, i.e.\ the empirical risk is a biased estimate of the true risk. A proof of this fact is found in Appendix~\ref{app: add res}. 
% non è vero ehhehee la devo ancora fare

We look for a bound that reproduces the shape of the empirical test error curve qualitatively. 
In particular, we want this function to have the same asymptotic behavior as $\expv_D[R(\hat h_D)]$,
as well as to be sensitive to changes in the structure of the learning problem which affect the test error curve. % For instance, as we will see, the test error changes with the noise level $(1-p)$ and with the complexity of the underlying data-generating function $h^*$, measured as the number of its sign changes in the interval $(0, 10$). 

All proofs are in Appendix~\ref{app: proofs}. 
\bigskip 


% we want to see if theory can explain the shape of the test error curve as model complexity increases, to explain generalization

% to do this, we want to find a function that bounds the test error and that reproduces its qualitative behavior 

% perché questo ci permette di capire quali aspetti del learning problem determinano quell'andamento

% in particular, cerchiamo un bound sul expected gap tra train e test error

\section{Union bound over $\cH$}\label{sec: unionbound}
In this section, we study the first bound on the quantity $\expv_D[(R-\hat R_D)(\hat h_D)]$, to which we refer as the {\it gap\/} for the rest of this work, without specifying that it is computed at $\hat h_D$ everytime.
This result is obtained starting from a concentration result (Hoeffding's inequality, see Theorem~\ref{app: hoeffding} in the Appendix) on the random variable $\hat R_D(h)$ for a fixed function $h\in\cH$. This is extended to a result on $\hat h_D$, which depends on the data $D$, by taking a union bound (see Proposition~\ref{app: unionbound} in the Appendix) on the probability that the event holds for any function $h\in\cH$.

\begin{theorem}\label{thm: unionbound}
\begin{equation}\label{eq: union bound}
    \expv_D[R(\hat h_D) - \hat R_D (\hat h_D)] \,\, \le \,\, O\left(\sqrt{\frac{\log \vert \mathcal H \vert}{n}}\right)
\end{equation}
\end{theorem}

{\bf Experimental design.} 
The preliminary experiments are discussed in detail and plots are provided in Appendix~\ref{app: prel exp}. 

We create the support $\mathcal X$ as a discretization over $N= 10,000$ numbers of the interval $[0,10]$. We sample from it with the uniform measure $p_\X(\bar x) = \sum_{\bar y\in \Y}p(\bar x, \bar y) = \frac{1}{N}$. 
The law to generate the labels is
\[
y \leftarrow \mathrm{sign}(h^*(x) + \sigma\varepsilon)
\] 
where $h^*:\X\to\Y$, is the data-generating function, $\varepsilon\sim\mathcal{N}(0,1)$, is a Gaussian noise, and $\sigma>0$. 
First, we generate the labels with a centered parabola, so that roughly half of them are $1$ and half $-1$ in the interval $[0,10]$. 
Since an algorithm could learn such a function with few datapoints (see Section~\ref{app/ sec: datagen}), we switch to a sinusoid with varying frequency,
\[
    h^*(x) = 5\, \sin\left(\frac{\pi \,( m  +1)}{10}\, x\right)
\]
where $m$ is the number of times $h^*$ crosses the $x$-axis in the open interval $(0,10)$, yielding $m+1$ contiguous blocks of $1s$ and $-1s$. 
We set as default value $m=10$.

We test a Bernoullian noise instead of a Gaussian one, concluding that there is no interesting difference in the performance of the algorithm as long as the portion of misclassified inputs due to the noise is the same. We find that for $m = 10$, $\sigma = 3.5$ achieves a noise of $20\%$, corresponding to $p =0.8$. For now, we stick to a gaussian noise.

In the end, our labels are generated as
 \[
    y \leftarrow \mathrm{sign}\left(5\, \sin\left(\frac{\pi \,( m  +1)}{10}\, x\right) + \sigma\varepsilon\right)
 \]
 with $\sigma = 3.5$ and $m = 10$. Also, our $\mathrm{sign}(\cdot)$ maps $0$ to $1$. 

We proceed as follows. We generate a test set with size $n_\text{test} = 200,000$ sampling from $\X$ with $p_\X$ and generating the labels as above. Then, for each depth $d \in [1, 7\lceil{\log_2 n}\rceil] \subset \mathbb N$ we do the following. For a number $T=100$ of iterations, we 
\begin{itemize}
    \item[$\cdot$] generate a train set ($n= 50$);
    \item[$\cdot$] fit a decision tree with depth at most $d$ (DecisionTreeClassifier algorithm from the scikit-learn library);
    \item[$\cdot$] compute the train error, the test error, and the differene between the two (the gap).
\end{itemize}
Then, we average the resulting quantities over the $T$ samples, obtaining empirical estimates of the expected values of the train error, test error, and gap over the disorder of the data. 

To plot the theoretical bound, instead, we need to express $\vert \mathcal H_d\vert$ as a function of $d$.
The exact computation gives
\[
\vert \mathcal H_d \vert = 2 \sum_{k=0}^{\min{N-1, 2^d-1}} \binom{N-1}{k}
\] 
\mati{Disclaimer: all the plots in this thesis have been obtained with $\vert\cH_d\vert \le 2^d$ which is incorrect, I need to fix it. I tried the formula above or something like $\min(2^N, N^{2^d}2^{2^d})$ but dependence on $d$ is awful.}

As the maximum depth $d$ increases, we plot the following:
\begin{enumerate}
    \item[(1)] the train and the test error curves;
    \item[(2)] the gap and the theoretical bound $\sqrt{\log \vert \mathcal H \vert/n}$;
    \item[(3)] the test error and the sum of train error and the bound (moving $\hat R_D(\hat h_D)$ to the right).
\end{enumerate}

\begin{figure}[h!]
  \centering
  \includegraphics[width=0.5\textwidth]{../../outputs/recap/pic1.png}
  \caption{$n = 50, N = 10,000, n_{\text{test}} = 200,000, \sigma = 5$, $h^*$ = sine with $10$ sign changes}
  \label{fig:pic1}
\end{figure}

We plot the bound~(\ref{eq: union bound}) to compare it with the empirical behavior of the test error in Figure~\ref{fig:pic1}. The first thing catching the eye is that the bound diverges, while the test error saturates. This divergence is due to the dependence of the bound on $\vert\cH_d\vert$, that is the number of binary trees with maximum depth $d$, which clearly increases with $d$. It follows immediately that a complexity measure of the hypotheses set that naively counts the functions it contains does not satisfy our requirements to find a qualitatively informative bound.

The test error of binary decision trees saturates because of a specific characteristic of the algorithm. As soon as the learned tree has only ``pure'' leaves, i.e.\ leaves containing inputs that have the same unique label in the training set, the algortihm ceases to split them and hence any additional allowed depth does not change the learned model. Clearly, the cardinality of the hypotheses sets $\cH_d, d>0$ does capture this property. 
% devo definire sopra cos'è H bene e H_d
Instead, a counting way that accounts for the finiteness of our problem might reproduce the saturation of the test error curve. Specifically, the configuration with only pure leaves can be achieved for two reasons: either all the possible values that the inputs can take are classified, or all the inputs in the sample are classified. Since, in our case, $n\ll N$, the second event happens much earlier than the first one. As a consequence, a complexity measure that counts the number of distinct classifications that the functions in $\cH$ can give, i.e.\ the vectors $h(x_1),\ldots,h(x_n) : h\in \cH$, instead of the functions themselves, could more faithfully represent the true complexity of the space of models that the algorithm can learn in practice.

\section{Rademacher Complexity and Massart Lemma}\label{sec: rademmassart}
In this section, we introduce the Rademacher complexity, a tool that allows to leverage the viewpoint exposed in the last section and obtain a saturating bound on the gap. Define $\LL =  \{\ell_h : h \in \mathcal H \}$ and $\ell_h(x,y) = \ell(h(x), y)$. 

\begin{definition}[Rademacher complexity]\label{def: rad}
Let $\xi_1,\ldots,\xi_n$ be Rademacher variables taking value $1$ or $-1$ with equal probability. The Rademacher complexity of the set $\LL$ is
\[
\mathcal R_n(\LL) = \expv_{\xi,D} \left[\sup_{\ell_h\in\LL}\frac{1}{n}\sum_{i=1}^n \xi_i \ell_h(X_i, Y_i)\right]
\]
\end{definition}

The Rademacher complexity measures the richness of a class of functions. The argument of the supremum  $\frac{1}{n}\sum_i\xi_i \ell_h(x_i,y_i) = \langle \xi, \ell_h(x,y)\rangle$ is the correlation between $\xi$ and $\ell_h(x,y)$. Since the loss function is fixed, this correlation measures to what extent the classifiers can align to random labels. Note that the supremum is taken before the expectation over the disorder of $\xi$ and $D$, and that in our case it is a maximum, since $\LL$ is finite. Hence, for each dataset and for each realization of the Rademacher vectors, we can choose a different function $\ell_h \in\LL$ that attains the maximum. A high complexity means that for most values of $\xi$ we can find a model $h$ that is similar to that pattern. It is the highest when we can find a function that correlates with any possible pattern.

The following results hold. 

\begin{theorem}[Rademacher complexity bound]\label{thm: rad bound}
%The gap between the expected and empirical risk of the ERM $\hat h$ is
\[
\expv_D[(R-\hat R) (\hat h)] 
% \le \expv_D[\sup_{h \in \cH}(R-\hat R) (h)] 
\le 2 \mathcal R_n(\LL)
\]
\end{theorem}

\begin{lemma}[Massart]\label{lemma: massart}
%We have
\[
\mathcal R_n(\LL) \le \expv_D\left[\sqrt{\frac{2\log N_D(\cH)}{n}}\right]
\]
where  $N_D(\cH)=\left|\left\{(\ell_h(X_1,Y_1),\ldots,\ell_h(X_n,Y_n)):h\in\cH\right\}\right|$.
\end{lemma}

As arises from these results, using the Rademacher complexity as an intermediate step allows us to replace the number of classifiers $\card$ with the number of loss vectors $N_D(\cH)$. In fact, the supremum in the Definition~\ref{def: rad} can be equivalently computed over the set of classifiers as well as over the set of loss vectors. The key to this result is that Theorem~\ref{thm: rad bound} bounds a quantity that depends on the true expected value of $\ell_h(X,Y)$ with a quantity that depends on loss vectors. This is possible thanks to a symmetrization technique, which consists in replacing the deviation of the random variable $\ell_h(X,Y)$ from its expected value with the difference between two random variables taking independent datasets $D$ and $D'$ as inputs. Note that since the loss function is fixed, counting loss vectors is equivalent to counting classification vectors. We work with loss vectors here because they arise more naturally from Rademacher complexity, and turn to classification vectors in Section~\ref{sec: growthf}. 

The following result encloses the key point of the saturation mechanism discussed at the end of Section~\ref{sec: unionbound}. 

\begin{lemma}[Counting inequalities]\label{lemma: counting}
\[
N_D(\cH) \le \min\{\card, 2^n\}.
\] 
\end{lemma} 

Substituting this bound on $N_D$ in Lemma~\ref{lemma: massart} and combining it with Theorem~\ref{thm: rad bound}, we obtain the following bound. 

\begin{theorem}[Rademacher and Massart bound]\label{thm: rademmassart bound}
\begin{align}
\expv_D[(R-\hat R) (\hat h_d)]&\le 2 \expv_D\left[\sqrt{\frac{2\log N_D(\cH)}{n}}\right] \notag \\
&\le 2\sqrt{\frac{2\min\{\log \vert\cH_d\vert, \log 2^n\}}{n}} \notag \\
&\le \min\left\{O\left(\sqrt\frac{\log\vert\cH_d\vert}{n}\right), O(1)\right\} \label{eq: rademmassart bound}
\end{align}
\end{theorem}

We set the constant term in bound~(\ref{eq: rademmassart bound}) to $0.5$, because the Rademacher complexity of $\LL$ is naturally bounded by $0.5$ (see Appendix~\ref{app: add res} for a proof of this fact). The rest of the experimental design is the same as before. 
\mati{si zia ma il bound è 2Rad quindi avrebbe più senso settarlo a $1$. però empiricamente il test error non va sopra 0.5. AAAAAA}

\begin{figure}[h!]
  \centering
  \begin{minipage}{0.48\textwidth}
    \centering
    \includegraphics[width=\textwidth]{../../outputs/recap/pic2.png}
    \caption{Identical parameters as Fig.1}
    \label{fig:pic2}
  \end{minipage}\hfill
  \begin{minipage}{0.48\textwidth}
    \centering
    \includegraphics[width=\textwidth]{../../outputs/recap/pic3.png}
    \caption{Different noise levels}
    \label{fig:pic3}
  \end{minipage}
\end{figure}

We plot the bound against the empirical error curve in Figure~\ref{fig:pic2}. In the leften part of the graph, the bound is the same as in Figure~\ref{fig:pic1}. From a depth of $13$ onwards, the minimum is instead attained by the constant term hence the bound saturates. This reflects the fact that beyond a certain depth, the number of distinct classification vectors no longer increases with $d$, through the mechanism explained above. 

At the same time, this mechanism is myopic to the other aspects of the structure of the learning problem.
Indeed, Figure~\ref{fig:pic3} shows how the gap curve is sensistive to different levels of noise in the dataset, as measured by the $\sigma$ coefficient (standard deviation). Specifically, it remains flatter with less noise, indicating both that the model performs better and that the difference between train and test error is lower (see the explorative experiments in Appendix~\ref{app: prel exp} showing how the empirical curves change with noise). The bound curve, instead, remains unchanged, failing to capture that the noise present in the data has an effect on the generalization performance of the algorithm. 

\section{Dependence on the Best Risk}\label{sec: bestrisk}
In this section, we implement a different approach and new theoretical tools to look for a noise-sensitive bound.

Let the minimum expected risk achievable in the hypotheses set be
\[
R^*_\cH = \min_{h\in\cH} R(h).
\]
We aim at obtaining a bound on the gap at $\hat h_D$ in terms of the best risk $R^*_\cH$. The first ingredient is a concentration result for $\hat R$ that depends on the true risk $R$.

\begin{lemma}[Relative deviation]
  For $\varepsilon>0$ with $n\varepsilon^2 \ge 8$,
  \[
  \Pr\left[\exists h\in\cH : (R-\hat R)(h) > \varepsilon\sqrt R(h)\right] \le 2N_{2n}(\cH)e^{-n\varepsilon^2/16}
  \]
  where $N_{n}(\cH) := \sup_{(x_1,y_1), \ldots,(x_{n},y_{n}) \in \X \times \Y} N_D(\cH)$.
\end{lemma}

The dependence on the supremum of the number of loss vectors $N_n$ emerges naturally, since the proof of this results uses a symmetrization technique.

In high probability, and isolating the empirical risk $\hat R$ on the right hand side, we obtain the following result. 

\begin{theorem}[Bound with empirical risk]\label{thm: bound emprisk}
    With probability at least $1-\delta$,
    \[
    R(\hat h) - \hat R(\hat h) \le \sqrt{\frac{\hat R(\hat h)\log(2N_{2n}(\cH)/\delta)}{n}} + 16\frac{\log(2N_{2n}(\cH)/\delta)}{n}
    \]
\end{theorem}

Now, let 
\[
h_0 \in argmin_{h\in\mathcal H} R(h)
\]
be the model that attains $R^*_\cH$. By definition of empirical risk minimizer, $\hat R_D(\hat h_D) \le \hat R_D(h_0)$. Applying Bernstein inequality to the variable $\hat R_D(h_0)$, we obtain a bound on $\hat R_D(\hat h_D)$ that depends on $\expv_D[\hat R_D(h_0)] = R^*_\cH$. Applying this result to the right hand side in the previous theorem, we obtain the desired bound depending on the best risk $R^*_\cH$. 

\begin{corollary}[Bound with best risk]\label{coroll: bound bestrisk}
    With probability at least $1-\delta$,
    \[
    R(\hat h) - \hat R(\hat h) \le \sqrt{\frac{R^*_{\cH}\log(4N_{2n}(\cH)/\delta)}{n}} + 20\frac{\log(4N_{2n}(\cH)/\delta)}{n}
    \]
\end{corollary}

Let the Bayes risk be
\[
R^* := \inf_{h:\X\to\Y} R(h)
\] 
where the infimum runs over a superset of $\cH$. The best risk can be decomposed as 
\[
R^*_\cH = \underbrace{R^*}_{\text{Bayes risk}} + \underbrace{R^*_\cH - R^*}_{\text{Approximation error}} =: \sigma^2 + A_\cH
\]
where
\begin{itemize}
\item the Bayes risk $\sigma^2$ measures the classification error due to the irreducible noise $\varepsilon$, as it is the minimal error that any function $h:\X\to\Y$ necessarily makes;
\item the approximation error $A_\cH$ is the excess error we make given that we consider a restricted set of classifiers, $\cH$.
\end{itemize}
Note that in the noise-less setting, $R^* = 0$ and the Bayes risk is attained by the data-generating function $h^*$. Moreover, if the classifier attaining $R^*$ lives in the hypotheses class, then $A_\cH = 0$. Thus, Corollary~\ref{coroll: bound bestrisk} gives a bound in terms of the irreducible noise and the approximation error. The next result states the final bound, making this dependence explicit and applying Lemma~\ref{lemma: counting} to garantee saturation.  

\begin{theorem}[Noise and approximation error bound]\label{thm: bound noiseandappr}
  \begin{equation}\label{eq: noiseandappr bound}
\expv_D\left[R(\hat h)-\hat R(\hat h)\right] \le \min\left\{ O\left(
\sqrt{\frac{(A_{\mathcal H}+\sigma^2)\log \card}{n}} + \frac{\log\card}{n} \right), O\left(
\sqrt{A_{\mathcal H}+\sigma^2}+1  \right)\right\}
\end{equation}
\end{theorem}

{\bf Experimental design:}
As compared to Sections~\ref{sec: unionbound} and~\ref{sec: rademmassart}, we turn to a label-flip noise with flip probability $1-p$. Therefore, we sample from $\X$ as before and we generate the labels as 
\[
y \leftarrow \sgn(h^*(x))b
\]
where $h^*$ is as above and $b$ is a random variable taking value $1$ with probability $p$ and $-1$ with probability $1-p$. We use $p = 0.8$ and $p\in\{0.6,0.8,0.95\}$ to inspect multiple levels of noise.

Since $p\ge 0.5$, the Bayes risk is 
\[
R^* = 1-p
\]
$R^*_\cH$ is computed exactly using dynamic programming. We produce labels using our sine function and the Bernoulli error and we cluster the intervals where labels have the same sign, obtaining $M$ contiguous blocks. We create a dp $k\times 2$ matrix that stores the minimum risk for each number of maximum splits $k$, up to $\min(M-1, 2^d)$ where $d$ is the maximum depth we are considering, when the last assigned label is $1$ or $-1$ (the columns). Then, for each contiguous block we dynamically update the dp matrix summing the minimum risk among all trees with at most $k$ splits. In the end, we select the minimum risk among the tree predicting $1$ or $-1$ in the very last block. 

The rest of the experimental design is as above.

\begin{figure}[h!]
  \centering
  \begin{minipage}{0.48\textwidth}
    \centering
    \includegraphics[width=\textwidth]{../../outputs/recap/pic4.png}
    \caption{Same setting, with Bernoullian noise at $20\%$}
    \label{fig:pic4}
  \end{minipage}\hfill
  \begin{minipage}{0.48\textwidth}
    \centering
    \includegraphics[width=\textwidth]{../../outputs/recap/pic5.png}
    \caption{The two terms of the minimum in the noise and approx. error bound}
    \label{fig:pic5}
  \end{minipage}
\end{figure}

\begin{figure}[h!]
  \centering
  \begin{minipage}{0.48\textwidth}
    \centering
    \includegraphics[width=\textwidth]{../../outputs/recap/pic6.png}
    \caption{Different noise levels}
    \label{fig:pic6}
  \end{minipage}\hfill
  \begin{minipage}{0.48\textwidth}
    \centering
    \includegraphics[width=\textwidth]{../../outputs/recap/pic7.png}
    \caption{Approximation error and Bayes risk}
    \label{fig:pic7}
  \end{minipage}
\end{figure}

We plot bound~(\ref{eq: noiseandappr bound}) in Figure~\ref{fig:pic4}. As before, the bound saturates due to the constant term. Figure~\ref{fig:pic5} shows the two terms in the minimum. Figure~\ref{fig:pic6} shows the bound and the empirical gap curve for different noise levels. Finally, we see that the bound changes with the noise. Both the term depending on $d$ and the constant one, which determines the saturating level, depend on the minimum risk $R^*_\cH$. As explained above, this changes with the noise, hence the bound capture the sensitivity of the generalization performance to different noise levels. 

\section{Growth function}\label{sec: growthf}

In this section, we introduce a new measure of the complexity of the hypotheses class, that considers classification vectors. This tool allows to make the intuition on the saturation of the test error due to characteristics of the decision tree algorithm that we discussed in Section~\ref{sec: rademmassart} more explicit. 

\begin{definition}[Growth function]
    The growth function of $\cH$ is
\[
 \Pi_{n}(\mathcal H):=\sup_{x_1,\ldots,x_n\in\mathcal X}
 \left|\left\{(h(x_1),\ldots,h(x_n)):h\in\cH\right\}\right|.
\]
\end{definition}

The following chain of inequalities extends Lemma~\ref{lemma: counting}. 
Given the set of images of a function $h$, applying a loss function to it can at most decrease the number of vectors in the set. At the same time, the images of the functions in $\cH$ are at most as many as the functions in $\cH$, and always less than or equal to $2^n$, the number of all possible $n$-long sequences of $1$ and $-1$.

The following extends Lemma~\ref{lemma: counting}.

\begin{lemma}[Counting inequalities with growth function]\label{lemma: growth}
For every sample $D$ of size $n$,
\[
 1\le N_D(\mathcal H)\le\Pi_{n}(\mathcal H)\le\min\{\vert \cH \vert,2^n\}.
\]
Moreover,
\[
\Pi_n(\mathcal H)  = N_n(\mathcal H) = \sup_{(x_1,y_1),\ldots,(x_n,y_n)\in \mathcal X \times \mathcal Y} N_D(\mathcal H)
\]
% qui D è la realizzazione forse, otherwise come la confronti con \Pi ?
\end{lemma}

The relation between $N_D$ and $\Pi_n$ makes it immediate to obtain growth-function versions of the bounds in Sections~\ref{sec: rademmassart} and~\ref{sec: bestrisk} (substitute $N(D)\le\Pi_n$ into Lemma~\ref{lemma: refined massart} and $N_{2n} = \Pi_{2n}$ into Theorem~\ref{thm: bound emprisk}). Substituting also $\Pi_{n}(\mathcal H)\le\min\{\vert \cH \vert,2^n\}$ yields bounds~(\ref{eq: rademmassart bound}) and (\ref{eq: noiseandappr bound}). The bounds are the same, but the intuition that they saturate because the number of distinct classification vectors is at most $2^n$ is made more explicit. In other words, when we evaluate the performance of an algorithm, the models in $\cH$ that the algorithm can learn are equivalent if they give the same classification. This result takes this into account when measuring the richness of the hypotheses class $\cH$.

\section{Going on}
The growth function has the additional result that
\[
 \log\Pi_{\cH}(2n)\le d\log\frac{2en}{d}
\]
we need to understand if this is useful. Plus, we mentioned the localized Rademacher sometime this summer. 

I still need to fix the matter of how to count $\vert \mathcal H \vert$ in terms of the depth $d$ (the ``disclaimer'' in Section~\ref{sec: unionbound}).
Setting the constant term at $0.5$ in bound~(\ref{eq: rademmassart bound}) is heuristical because I know it works. Needs motivation.
Similarly, I say $O(\sqrt{A_{\mathcal H}+\sigma^2} + 1) = O(\sqrt{A_{\mathcal H}+\sigma^2})$ and plot only $\sqrt{A_{\mathcal H}+\sigma^2}$ but it's heuristical as well, and since $A_{\mathcal H}+\sigma^2 \le 1$ probably wrong (however, con il $+1$ non fa nemmeno in tempo a saturare perché l'altro termine che dipende da $d$ rimane più piccolo). 

\newpage
\appendix
\section{Preliminary experiments}\label{app: prel exp}

\subsection{Noise structure}
We generate the labels with a centered parabola (see Section~\ref{app/ sec: datagen}), so that roughly half of them are $1$ and half $-1$ in the interval $[0,10]$. Then, we inspect what happens if we use a Bernoullian noise instead of a Gaussian one, i.e. when data is generated as 
\[
y \leftarrow sign(h^*(x))b
\]
for $b\sim\mathrm{Bernoulli}(p)$. Notice that this noise flips the labels with a probability of $1-p$. 
While from the first experiments it seemed that the algorithm achieved a higher test error, we postulate that this may be due to the portion of misclassified inputs, rather than to the noise structure itself. Hence, fixing $p = 0.8$, we look for the value of $\sigma$ that achieves roughly the same portion of mislcassified inputs ($20\%$). Figure~\ref{app/fig: sigmastar parabola} shows that this happens for $\sigma = 5.5$, and Figures~\ref{app/fig: parabola bernoulli} and~\ref{app/fig: parabola gauss} show that the two noise structures then yield essentially the same train and test errors. 

\begin{figure}[h!]
  \centering
  \includegraphics[width=0.8\textwidth]{../../outputs/recap/app_sigmastar_parabola.png}
  \caption{Average fraction of misclassified inputs against $\sigma$, over $100$ samples of size $n = 50$, $h^*$ = centered parabola}
  \label{app/fig: sigmastar parabola}
\end{figure}

\begin{figure}[h!]
  \centering
  \begin{minipage}{0.48\textwidth}
    \centering
    \includegraphics[width=\textwidth]{../../outputs/recap/app_parabola_bernoulli.png}
    \caption{$n = 50, N = 10,000, n_{\text{test}} = 200,000$, Bernoulli noise with $p = 0.8$}
    \label{app/fig: parabola bernoulli}
  \end{minipage}\hfill
  \begin{minipage}{0.48\textwidth}
    \centering
    \includegraphics[width=\textwidth]{../../outputs/recap/app_parabola_gauss.png}
    \caption{$n = 50, N = 10,000, n_{\text{test}} = 200,000$, Gaussian noise with $\sigma = 5.5$}
    \label{app/fig: parabola gauss}
  \end{minipage}
\end{figure}

\subsection{Data-generating function}\label{app/ sec: datagen}
Our first choice of $h^*$ is a parabola centered horizontally at $5$, the midpoint of the support $\mathcal X$, and vertically so that roughly half of the points fall above the $x$-axis and half below. I.e.,
\[
h^*(x) = (x-5)^2 -6.25
\]

\begin{figure}[h!]
  \centering
  \includegraphics[width=0.7\textwidth]{../../outputs/recap/app_explore_parabola.png}
  \caption{A sample of size $n = 50$ generated with the centered parabola, for $\sigma \in \{10^{-4}, 1, 5.5, 10\}$}
  \label{app/fig: explore data parabola}
\end{figure}

Figure~\ref{app/fig: explore data parabola} shows the data generated with the centered parabola, for different amounts of noise.

Because in theory an algorithm could learn such a classification pattern as soon as it sees $3$ points (one point for each region), we investigate what happens when the underlying function is more complex. We study a sinusoid with varying frequency,
\[
    h^*(x) = 5\, \sin\left(\frac{\pi \,( m  +1)}{10}\, x\right)
\]
We multiply the sine by $5$ to allow for enough height to properly inspect the data, while the argument of the sine is such that $m$ is the number of times $h^*$ crosses the $x$-axis in the open interval $(0,10)$, yielding $m+1$ contiguous blocks of $1s$ and $-1s$. 

\begin{figure}[h!]
  \centering
  \begin{minipage}{0.32\textwidth}
    \centering
    \includegraphics[width=\textwidth]{../../outputs/recap/app_sine_m10.png}
    \caption{$m = 10$}
    \label{app/fig: sine m10}
  \end{minipage}\hfill
  \begin{minipage}{0.32\textwidth}
    \centering
    \includegraphics[width=\textwidth]{../../outputs/recap/app_sine_m100.png}
    \caption{$m = 100$}
    \label{app/fig: sine m100}
  \end{minipage}\hfill
  \begin{minipage}{0.32\textwidth}
    \centering
    \includegraphics[width=\textwidth]{../../outputs/recap/app_sine_m1000.png}
    \caption{$m = 1000$}
    \label{app/fig: sine m1000}
  \end{minipage}
\end{figure}

\begin{figure}[h!]
  \centering
  \includegraphics[width=0.8\textwidth]{../../outputs/recap/app_ratio_sweep.png}
  \caption{$n = 50, N = 10,000, n_{\text{test}} = 200,000$, $\sigma = 3$, train and test error for different ratios $r = m/n$}
  \label{app/fig: ratio sweep}
\end{figure}

We explore how a random sample and the performance of the algorithm change for different sine frequencies. In particular, we expect that for a high frequency the algorithm will not learn, as our datasize is $n=50$.
In Figures~\ref{app/fig: sine m10}--\ref{app/fig: sine m1000}, we test a function that osciallets $10, 100$, and $1000$ times with a noise level of $\sigma = 3$ (we reduce it compared to $\sigma = 5$, having increased the complexity of the function). We notice that for $m=1000$ the algorithm fails to learn, and the test error stays stable at $0.5$. This is because the labels change so frequently with respect to the number of points in the training dataset $n$ that it is as if we had random labels. Hence, we investigate the relation between the complexity of the underlying function $m$ and the data size $n$. We find that for $\frac{m}{n}>5$, roughly, the algorithm ceases to learn (Figure~\ref{app/fig: ratio sweep}). In the following experiments, we fix $m=10$ and $n=50$, which is very far from the threshold.

\begin{figure}[h!]
  \centering
  \includegraphics[width=0.8\textwidth]{../../outputs/recap/app_misclassified_sine.png}
  \caption{Fraction of misclassified inputs over $100$ samples of size $n = 50$, $h^*$ = sine with $m = 10$; dashed lines are the averages}
  \label{app/fig: misclassified sine}
\end{figure}

\begin{figure}[h!]
  \centering
  \includegraphics[width=0.8\textwidth]{../../outputs/recap/app_sigmastar_sine.png}
  \caption{Average fraction of misclassified inputs against $\sigma$, over $100$ samples of size $n = 50$, $h^*$ = sine with $m = 10$}
  \label{app/fig: sigmastar sine}
\end{figure}

We then verify whether a Gaussian noise with $\sigma = 5.5$ achieves the same portion of misclassified inputs on average as a Bernoullian noise with $p=0.8$, as was the case for data generated with a parabola. Figure~\ref{app/fig: misclassified sine} shows that this is not the case on average ($29\%$ against $20\%$), while Figure~\ref{app/fig: sigmastar sine} shows that the desired value is $\sigma = 3.5$.

\begin{figure}[h!]
  \centering
  \includegraphics[width=0.8\textwidth]{../../outputs/recap/app_simple_function_sigma1.png}
  \caption{$N = 10,000, n_{\text{test}} = 200,000, \sigma = 1$, $h^*$ = sine with $m = 2$: test error at depth $2$ and $3$ against $n$, averaged over $100$ samples}
  \label{app/fig: simple function sigma1}
\end{figure}

\begin{figure}[h!]
  \centering
  \includegraphics[width=0.8\textwidth]{../../outputs/recap/app_simple_function.png}
  \caption{$N = 10,000, n_{\text{test}} = 200,000, \sigma = 3.5$, $h^*$ = sine with $m = 2$: test error at depth $2$ and $3$ against $n$, averaged over $100$ samples}
  \label{app/fig: simple function}
\end{figure}

Finally, we investigate the afore-mentioned intuition about how easy it is for the algorithm to learn a simple function. We take $m=2$ (which corresponds to a parabola, in terms of classification pattern), sample uniformly from $\X$, and look for the training set size $n$ for which the test error gets within $0.01$ of the Bayes risk on average, at depth $2$ and $3$.

We start with little noise, $\sigma = 1$, for which the Bayes risk is $\approx 0.05$ (Figure~\ref{app/fig: simple function sigma1}). Here the algorithm learns the function quickly: with $n = 50$ points the test error is already $\approx 0.07$, within $0.02$ of the Bayes risk, and it gets within $0.01$ with $n = 200$ at depth $2$ and $n = 500$ at depth $3$.

With our noise level, $\sigma = 3.5$, the Bayes risk is $\approx 0.20$ and $n = 1000$ training points are needed on average for the test error to get within $0.01$ of it, both at depth $2$ and $3$ (Figure~\ref{app/fig: simple function}). Hence, although two splits are enough to represent the function, with this amount of noise the algorithm needs many points to place them correctly.


\newpage
\section{Proofs}\label{app: proofs}
\begin{proof}[Proof of Theorem~\ref{thm: unionbound}]
  Apply Hoeffding Lemma and a union bound over the hypotheses set $\mathcal H$. Tail integrate to obtain a result in expected value. 
\end{proof}

\begin{proof}[Proof of Theorem~\ref{thm: rad bound}]
  Symmetrization technique. 
\end{proof}

\begin{proof}[Proof of Lemma~\ref{lemma: refined massart}]
  Pass to the supremum over {\it loss vectors\/} instead of functions $\ell\circ h$, and bound it with the sum, which makes the $N(D)$ term appear. Then, apply the exponential Hoeffding Lemma and optimize for the coefficient in the exponent. 
\end{proof}

\begin{proof}[Proof of Lemma~\ref{lemma: counting}]
  Given a dataset, its realizations through $\ell_h$ are less than or equal to those through $h$ because applying $\ell$ cannot increase the number of outputs, and the number of images of the functions in $\mathcal H$ are less than or equal to the number of functions $\card$. On the other hand, $N_D(\cH)\le 2^n$ because we have at most $2^n$ $n$-long sequences of $0s$ and $1s$.
\end{proof}

\begin{proof}[Proof of Theorem~\ref{thm: rademmassart bound}]
  Apply Theorem~\ref{thm: rad bound}, Lemma~\ref{lemma: refined massart} and Lemma~\ref{lemma: counting}. 
\end{proof}

\begin{proof}[Proof of Theorem~\ref{thm: bound emprisk}]
        Apply the relative deviation inequality (see Loz's file),
        \[
        \Pr\left[\exists h\in\cH : R(h)-\hat R(h) > \varepsilon \sqrt{R(h)}\right] \le 2N_{2n} (\cH) e^{-2\varepsilon^2/16}
        \]
        with the suitable $\varepsilon$ and then isolate $\hat R(\hat h)$ on the RHS.
        % I have not looked precisely at the proof of the rel dev ineq
\end{proof}

\begin{proof}[Proof of Corollary~\ref{coroll: bound bestrisk}]
         Set $h_0 = \argmin_{h\in\cH} R(h)$, so that $R(h_0) = R^*_{\cH}$. Then,
        \begin{equation}\label{eq: hatR h le hat R h0}
        \hat R(\hat h) \le \hat R(h_0).
        \end{equation}
        Since $h_0$ is independent of the data, $\expv_D[\hat R(h_0)] = R(h_0)$ and hence Bernstein gives, with probability at least $1-\delta/2$,
        \[
        \hat R(h_0) \le R^*_{\cH} + \sqrt{\frac{2R*_\cH \log (2/\delta)}{n}} + \frac{\log(2/\delta)}{3n}.
        \]
        Since $N_{2n}(\cH)\ge 1$ and thus $\frac{\log (2/\delta)}{n}\le \frac{\log(4N_{2n}(\cH)/\delta)}{n}$,
        \begin{equation}\label{eq: bound hatR h0}
        \hat R(h_0) \le \left(\sqrt{R^*_\cH} + \sqrt{\frac{4\log(N_{2n}(\cH)/\delta)}{n}}\right)^2
        \end{equation}
        Substitute~(\ref{eq: hatR h le hat R h0}) into~(\ref{eq: bound hatR h0}), and apply the bound to $\hat R(\hat h)$ in the RHS of Theorem~\ref{thm: bound emprisk} at confidence $1-\delta/2$. This yields the desired bound, which holds with probability at least $1-\delta$ by a union bound on the events ((\ref{eq: bound hatR h0}) and Theorem~\ref{thm: bound emprisk}) that we require to hold simultaneously.
\end{proof}

\begin{proof}[Proof of Corollary~\ref{coroll: bound noiseandappr}]
  Substitute $R^*_\cH = \sigma^2 + A_\cH$ into Corollary~\ref{coroll: bound bestrisk}.
\end{proof}

\begin{proof}[Proof of Theorem~\ref{thm: bound noiseandappr}]
    Let $\text{gap} = (R-\hat R(\hat h))$ and $C = \sqrt{\frac{\log(4 N_n(\cH)/\delta)}{n}}$. Corollary~\ref{cor: best risk} gives
\begin{align*}
&\Pr\left[\text{gap} \le \sqrt{R^*_\cH}\,C + 20C^2\right] \ge 1-\delta \\
\implies &\Pr\left[\text{gap} > \sqrt{R^*_\cH}\,C + 20C^2\right] \le \delta
\end{align*}
Set $\varepsilon = \sqrt{R^*_\cH}\,C + 20C^2$. Let us now find $\delta(\varepsilon)$:
\[
20C^2 + \sqrt{R^*_\cH}\,C - \varepsilon = 0 \implies C = \frac{-\sqrt{R^*_\cH} + \sqrt{R^*_\cH + 80\varepsilon}}{40}
\]
So we have
\begin{align*}
\frac{\log(4N_{2n}(\cH)/\delta)}{n} &= \left(\frac{\sqrt{R^*_\cH + 80\varepsilon} - \sqrt{R^*_\cH}}{40}\right)^2 \\
\frac{4N_{2n}(\cH)}{\delta} &= \exp\left(n\left(\frac{\sqrt{R^*_\cH + 80\varepsilon} - \sqrt{R^*_\cH}}{40}\right)^2\right)
\end{align*}
and hence
\[
\delta(\varepsilon) = 4N_{2n}(\cH)\exp\left(-\frac{n}{1600}\left(\sqrt{R^*_\cH + 80\varepsilon} - \sqrt{R^*_\cH}\right)^2\right).
\]
Using that $\expv[X] = \int_0^\infty \Pr[X \ge \varepsilon]\,d\varepsilon$ for $X \ge 0$, applied to $\text{gap}^+ \coloneq \max\{\text{gap}, 0\}$,
\[
\expv_D[\text{gap}] \le \expv_D[\text{gap}^+] = \int_0^\infty \Pr[\text{gap} \ge \varepsilon]\,d\varepsilon \le \int_0^\infty \min\{1, \delta(\varepsilon)\}\,d\varepsilon.
\]
By a shift of perspective, looking at the integral as a sum of horizontal stripes instead of vertical stripes in the graph of $\varepsilon \mapsto \delta(\varepsilon)$, we obtain
\[
\int_0^\infty \min\{1, \delta(\varepsilon)\}\,d\varepsilon = \int_0^1 \delta^{-1}(t)\,dt = \int_0^1 \varepsilon(\delta)\,d\delta
\]
Hence, writing $\varepsilon(\delta) = \sqrt{R^*_\cH}\,C + 20C^2$ explicitly,
\[
\expv_D[\text{gap}] \le \int_0^1 \varepsilon(\delta)\,d\delta = \sqrt{\frac{R^*_\cH}{n}}\int_0^1 \sqrt{\log\frac{4N_{2n}(\cH)}{\delta}}\,d\delta + \frac{20}{n}\int_0^1 \log\frac{4N_{2n}(\cH)}{\delta}\,d\delta.
\]
The logarithm is the only part depending on $\delta$, and its integral is
\[
\int_0^1 \log\frac{4N_{2n}(\cH)}{\delta}\,d\delta = \log(4N_{2n}(\cH)) - \int_0^1 \log\delta\,d\delta = \log(4N_{2n}(\cH)) + 1.
\]
Applying Jensen's inequality to the square root term,
\[
\int_0^1 \sqrt{\log\frac{4N_{2n}(\cH)}{\delta}}\,d\delta \le \sqrt{\int_0^1 \log\frac{4N_{2n}(\cH)}{\delta}\,d\delta} = \sqrt{\log(4N_{2n}(\cH)) + 1},
\]
we obtain
\[
\expv_D[\text{gap}] \le \sqrt{\frac{R^*_\cH\left(\log(4N_{2n}(\cH)) + 1\right)}{n}} + 20\,\frac{\log(4N_{2n}(\cH)) + 1}{n}.
\]

Asymptotically,
\[
\expv_D\left[R(\hat h)-\hat R(\hat h)\right] \le O\left(
\sqrt{\frac{(A_{\mathcal H}+\sigma^2)\,\log N_{2n}(\mathcal H)}{n}} + \frac{\log N_{2n}(\mathcal H)}{n} \right)
\]
because, at the numerator,
\[
\log (4 N_{2n}(\mathcal H)) + 1 = \log 4 + \log N_{2n}(\mathcal H) + 1 = \log N_{2n}(\mathcal H) + O(1)
\]
and for large $x$, 
\[
\log x + O(1) \le C \log x
\] 
for some constant $C>0$. Plugging in the counting inequality from Lemma~\ref{lemma: counting} with $2n$, we obtain and plot the following bound
\[
\expv_D\left[R(\hat h)-\hat R(\hat h)\right] \le \min\left\{ O\left(
\sqrt{\frac{(A_{\mathcal H}+\sigma^2)\log \card}{n}} + \frac{\log\card}{n} \right), O\left(
\sqrt{A_{\mathcal H}+\sigma^2}+1  \right)\right\}
\]
Remark: recall that $N_{2n}(\cH)$ is the maximum number of loss vectors over the union of two datasets $D, D'$, or equivalently over a dataset with size $2n$. The counting lemma must be applied to $2n$-long loss vectors. Hence, we still have $N_{2n}(\cH) \le \vert \cH\vert$ because the number of functions is still the same, and the bound counting all the possible sequences becomes $N_{2n}(\cH) \le 2^{2n}$. 

\end{proof}

\begin{proof}[Proof of Lemma~\ref{lemma: growth}]
    $N_D(\cH) \le \Pi_{n}(\cH)$ because the number of distinct loss vector can at most decrease, given prediction vectors. $\Pi_{n}(\cH)\le \card$ because the realizations of functions in $h$ are at most as much as the functions in $h$ (equality when functions in $h$ are bijective). $\Pi_{n}(\cH)\le 2^n$ because we have at most $2^n$ sequences of $n$ $1$s and $-1$s.
\end{proof}

\newpage
\section{Additional results}\label{app: add res}

\begin{fact}[The empirical risk is a biased estimate of the true risk for the ERM]
lalala
  
\end{fact}

\begin{fact}[Bound on the Rademacher complexity of $\LL$]
lalalala

\end{fact}


\end{document}